
% ============================================================================
\clearpage
\section*{v4.10.0-r14 Source-Only: Status Notes on the Koide Reference, Bent Axes and the Guarded Core}
\addcontentsline{toc}{section}{v4.10.0-r14 Source-Only: Status Notes on the Koide Reference, Bent Axes and the Guarded Core}
\label{sec:v4100-r14-bent-axes-guarded-core}
% ============================================================================

\paragraph{Purpose.}
After v4.10.0-r13 the OP3 work continued from the profile that the present inventory could not meet.  This layer records
the explorations of 14 and 15 November 2026: the slot Fourier operation and its limits, the search for a source of the
balanced reference, a first map between the sedenion and the Dixon carriers, the reading of the charged-lepton phase as a
bent octant, and the reading of the bend as a flux-free guarded core in the Hopf fibration.  It records one test (H104,
refuted) and two corrections of statements made during the work.  Apart from H104 all items are explorations without
pre-registration; none is compared with data, and the Koide target \(K=\tfrac23\) and the phases were known in advance.  The
notes are \texttt{notes/\allowbreak research/\allowbreak v4\_10\_0\_r14/} and the folders and files cited there.  Earlier layers
are not rewritten.

\noindent\textbf{Source-only status.}
This v4.10.0-r14 layer is source-only.  It does not build a PDF, does not assign a new DOI,
and does not authorize a public release.  The v5.0.0 Version DOI
\href{https://doi.org/10.5281/zenodo.24152615}{10.5281/zenodo.24152615} and the Concept DOI
\href{https://doi.org/10.5281/zenodo.19210728}{10.5281/zenodo.19210728} are unchanged by it.

\subsection*{v4.10.0-r14 status notes}
\label{subsec:v4100-r14-status}

\paragraph{Theoremlet v4.10.0-r14.1 (H104: the singular position of the counting pair depends on the fixed spectrum).}
Pre-registered (commit \texttt{c58f449}).  With a free area spectrum (23 rules, 253 pairs) a second pair of counting rules
meets all three windows of v4.10.0-r09.1 (\(\gamma_L=1.864824\) against \(1.864700\)).  H104 is refuted.  H98 remains confirmed
in its own rule space; the singular position of the pair rests on the fixed spectrum.  A first run with a faulty solver had
reported ``confirmed''; it is disclosed and kept.

\paragraph{Theoremlet v4.10.0-r14.2 (the slot Fourier operation and its limits).}
On the slot space \(\mathbb C^3\) let \(\tau\) be the cyclic shift and \(D=\operatorname{diag}(1,\omega,\omega^2)\).  The Fourier
matrix \(F\) of the three slots is the equally weighted sum of the nine operations \(\tau^aD^b\); \(F^2\) is the slot part of
\(\rho\).  With the phase \(S=\operatorname{diag}(1,\omega,\omega)\), \(F\) and \(S\) generate \(2T\) on the doublet, the axes meet
at \(|\cos|=1/\sqrt3=\langle e_0|F|e_0\rangle\), and the \(S\)-orbit of the Fourier-fixed vector \((1+\sqrt3,1,1)\) has
\(K=\tfrac23\) and Bargmann phase \(\pi/4\).  \(F\) and \(S\) are assumed as operations.  \(F\) does not preserve the triality
form, \(\tau\)-preserving dynamics cannot produce \(F\), and the static two-probe channel admits no return of \(F\).  The
reference is selected as the unique ground state of \(H=-(aQ+bR)\), \(Q=(1+D+D^2)/3\), \(R=(1+\tau+\tau^2)/3\), with
\(K_{\rm orb}=\tfrac23\) exactly for \(a=b\); a fixed \(H\) breaks \(S\), and the orbit Gram matrix has rank two.

\paragraph{Theoremlet v4.10.0-r14.3 (no source for the balance \(a=b\)).}
No action in the inventory supplies \(Q\) and \(R\) with bound equal couplings.
\begin{enumerate}[leftmargin=2em]
\item[(a)] Monomial slot symmetries (1296 elements with sixth roots of unity) map \(\langle D\rangle\) only to itself and cannot
exchange \(Q\) and \(R\).
\item[(b)] Every equally weighted pair of eigenlines from two different cyclic subgroups of the Weyl group (mutually
unbiased bases) selects the octant reference; all 54 such pairs do.  Only the equal weight is open, and it is the
\(F\)-self-duality of a \(\mathbb Z_3\) clock term \(-J(Z+Z^\dagger)-h(X+X^\dagger)\) at \(J=h\).
\item[(c)] The mixing mask of v4.9.0-r47 has the same form.  In the exact forms of v4.9.0-r52 the rank-one reference
needs \(\varepsilon^*=\tfrac13\ln\frac{1+\sqrt3}{2}\approx0.104\), whereas the mixing fit of v4.9.0-r54 gives
\(\varepsilon_e\in[0.44,0.82]\).
\item[(d)] Frame readings: with \(a\), \(b\) as frame weights the balance is the photon weight; the frame maps \(R\), \(T\),
\(J_B\) close in \(\mathfrak{sl}(2,\mathbb R)\), and their ordered products are \(\pm1\) without a phase.  The assumption of
v4.9.0-r33 is not the bottleneck: every frame plane is commutative, and the frame triole is a transposition (v4.9.0-r38).
The native Levi-Civita triple with a Hurwitz cycle is \(\operatorname{Im}\mathbb H\) of v4.9.0-r33.1, which breaks \(Q_{\rm em}\);
under the charge axis \(e_8\) the neutral triple \(R_{e_7},R_{e_8},R_{e_{15}}\) carries two generations only (r13.4).
\end{enumerate}

\paragraph{Theoremlet v4.10.0-r14.4 (a colour- and charge-preserving map \(\mathbb C\otimes\mathbb S\to\mathbb C\otimes\mathbb H\otimes\OO\)).}
Under CT-003, \(\mathbb C\otimes\mathbb S\) carries \(\{\nu,e^+,u,\bar d\}\) twice.  Injective linear maps that intertwine
\(\mathfrak{su}(3)_c\) and send \(Q_*\) to \(Q_{\rm em}\) form a 20-dimensional family.  The naive Cayley--Dickson embedding
preserves colour but not charge.  The neutral family \(\mathfrak{su}(2)\) of r13.4 maps onto the left action of \(\mathbb H\) (the
Lorentz spinor of v4.9.0-r33.2) sector by sector, with opposite orientation on leptons and quarks, and never onto the right
action (weak isospin) or the conjugation of v4.9.0-r33.1(b).  Reading: compared through colour and charge, the two
generations of route A occupy the place of the spin doublet; colour and charge alone do not separate generation and spin.

\paragraph{Theoremlet v4.10.0-r14.5 (bent axes).}
Three states with pairwise overlap \(|\langle a|b\rangle|^2=\tfrac12\) give \(K=\operatorname{Tr}G^2/(\operatorname{Tr}G)^2=\tfrac23\)
for every Bargmann phase \(\varphi\), and
\[
\det G=\tfrac12\bigl(\sqrt2\cos\varphi-1\bigr).
\]
At \(\varphi=\pi/4\) (the flat octant) \(\det G=0\): the states lie in one qubit and one mass vanishes.  For \(\varphi<\pi/4\)
they span \(\mathbb C^3\) and give three positive masses; \(\varphi=\tfrac23\) reproduces the charged-lepton ratios
\(1:14.38:58.97\).  The bend \(\pi/4-\tfrac23=0.118731\) is the flux of the placeholders CT-001/CT-002, read as
\(\gamma_M/2=0.118766\).  On the Bloch sphere three points at \(90^\circ\) form a rigid octant; a smaller enclosed area requires
leaving the sphere for \(\mathbb C^3\).

\paragraph{Theoremlet v4.10.0-r14.6 (the guarded core as a flux-free cap).}
Read the guarded core of the orbit scan (\(X=S^3\setminus B_\varepsilon(K)\), proposition ``Hopf completion versus
guarded-core remainder'') as a cap of solid angle \(\Omega\) on the Bloch sphere of the generations that contains no flux.
\begin{enumerate}[leftmargin=2em]
\item[(a)] With the core inside the generation triangle the phase falls below \(\pi/4\), with the core outside it does not;
the sign follows from the position, and inside the triangle the position does not matter.
\item[(b)] Topology: on the sphere without \(n\) caps there are \(n-1\) free Aharonov--Bohm holonomies.  No core gives no bend,
two or more cores leave the phase free; exactly one core is the only number that bends without a free parameter.
\item[(c)] With one core the displaced flux stays on the sphere.  Spread evenly it gives \(\varphi=0.679475\); placed outside the
triangle, for instance at the antipode, it gives \(\varphi=\pi/4-\Omega/2\), the relation of H102 for \(\Omega=\gamma_M\).  The
two curves are proportional with factor \(4\pi/(4\pi-\gamma_M)\); core size and the fraction placed outside are degenerate,
and \(\Omega=\gamma_M\) holds only together with full placement outside.
\item[(d)] The antipode follows from the reflector.  By v4.9.0-r02.2 and r02.3 the involutive Umkehrwalze is antilinear and
fixes no field.  On a spin-\(\tfrac12\) sphere an antilinear map \(\Theta\) with \(\Theta^2=-1\) is the antipodal map, without fixed
points; with \(\Theta^2=+1\) it is a reflection with fixed points.  Read as a reflector, the core returns the flux at the
antipode, point to point and not spread.  This is a reading.
\item[(e)] The ball \(B_\varepsilon(K)\) removes no fibre for \(\varepsilon<\pi\) but breaks the fibres over a cap of radius
\(2\varepsilon\).  \(\Omega=\gamma_M\) corresponds to \(\varepsilon=7.90^\circ\) (a cut fibre has no holonomy) or to a volume fraction
\(\gamma_M/(4\pi)\) of \(S^3\) (proportional reading, \(\varepsilon=25.94^\circ\)).
\end{enumerate}

\paragraph{Theoremlet v4.10.0-r14.7 (the core size is OP1).}
The counting rule \(\sum_j2e^{-2\pi\gamma\sqrt{j(j+1)}}=1\) gives \(\gamma_M\) as an opening angle in the plane of a puncture,
not as a solid angle on the generation sphere.  The geometry of quantum isolated horizons (Ashtekar, Baez, Krasnov,
arXiv:gr-qc/0005126; Engle, Noui, Perez, arXiv:0905.3168) has the same structure, a punctured sphere with a U(1)
connection whose curvature sits in the punctures, but fixes no angle \(\gamma_M\): the defect \(4\pi a/k\) depends on the horizon
area, and \(k=4\pi/\gamma_M=52.9\) is not an integer.  Reading the core as the area quantum of the observer frame gives
\(\gamma_M\) for every mass measured in the matter frame, including neutrinos, which the main body reads there.  The core size
thereby coincides with OP1: why the matter frame counts all spins and the tachyon frame only the floor.

\paragraph{Corrections.}
\begin{enumerate}[leftmargin=2em]
\item During the work it was stated that H99 places the neutrinos on the octant.  That is wrong: H99 has
\(\varphi_\nu=\varphi_e+\pi/4\) with a negative root.  The neutrino figure is two octants minus the same core; the idea that charge
places the core inside the triangle is dropped.
\item It was first stated that with a removed core the flux leaves over the boundary.  With exactly one core the remainder is
simply connected and the flux stays on the sphere; r14.6(c) is the corrected form.
\end{enumerate}

\subsection*{v4.10.0-r14 status in the three frames}
\label{subsec:v4100-r14-frames}

\begin{itemize}[leftmargin=2em]
\item M: all mass statements are matter-frame readouts; r14.7 reads the core size as the matter-frame quantum.
\item T: \(\gamma_T=\gamma_M/\gamma_L\) is the tachyon-frame quantum; no statement is transported to T.
\item P: the balance of r14.3(d) is the photon weight; otherwise not evaluated.
\end{itemize}

\subsection*{v4.10.0-r14 register entries}
\label{subsec:v4100-r14-register}

\begin{center}
\small
\begin{tabular}{@{}>{\raggedright\arraybackslash}p{0.34\textwidth}>{\raggedright\arraybackslash}p{0.58\textwidth}@{}}
\toprule
statement & status after v4.10.0-r14\\
\midrule
H104 & refuted (r14.1)\\
slot Fourier \(F\) and reference selection & construction with assumed operations; balance \(a=b\) without source (r14.2, r14.3)\\
map \(\mathbb C\otimes\mathbb S\to\mathbb C\otimes\mathbb H\otimes\OO\) & exists, not unique; family \(\mathfrak{su}(2)\) at the place of spin (r14.4)\\
bent axes & exact: \(\det G=\tfrac12(\sqrt2\cos\varphi-1)\) (r14.5)\\
guarded core & reading: flux-free, one core, displaced flux at the antipode (r14.6)\\
core size \(\gamma_M\) & open; coincides with OP1 (r14.7)\\
CT-001, CT-002, CT-003 & unchanged; none structurally matched\\
H100 / H102 & pending / pending\\
\bottomrule
\end{tabular}
\end{center}

\subsection*{v4.10.0-r14 ledger and guard update}
\label{subsec:v4100-r14-ledger-guard}

\[
R_{\rm chk}^{v4100\text{-}r14}=(1,1,0)
\]
for (status notes recorded; formulas checked in the cited notes; derivation of the overlap \(\tfrac12\), of the core size and
of a three-generation carrier open).

The layer spends no guard:
\[
R_{\rm glob}^{(v4100\text{-}r14,+)}=(0,1,1,1,1),
\qquad
\Gamma_{\rm guard}^{v4100\text{-}r14}=(1,1,1,1,1),
\qquad
\beta_4=1/256\ \text{unspent}.
\]
OP~1, OP~3 and OP~4 remain open.  The layer does not imply exact finite \(C5\) zero closure, exact all-readout terminality,
\(C6/T2\), a global photon-clock/RH theorem, or global all-readout closure.
